\chapter{TỔNG QUAN ĐỀ TÀI VÀ BỐI CẢNH ỨNG DỤNG}

\section{Bối cảnh thực tiễn}

Trong kỷ nguyên số hóa mạnh mẽ, hoạt động tiếp thị (marketing) của doanh nghiệp đã chuyển dịch từ các chiến dịch đơn lẻ sang quy trình vận hành đa kênh (omni-channel) liên tục. Một chiến dịch marketing điển hình đòi hỏi sự phối hợp nhịp nhàng giữa nhiều mắt xích: hoạch định mục tiêu kinh doanh, phân bổ ngân sách tài chính, lựa chọn sản phẩm trọng tâm và 7 nhóm kênh tiếp cận (\texttt{FACEBOOK}, \texttt{INSTAGRAM}, \texttt{TIKTOK}, \texttt{ZALO\_OA}, \texttt{EMAIL}, \texttt{GOOGLE\_ADS}, \texttt{WEBSITE}), sáng tạo thông điệp, kiểm duyệt nội dung, lập lịch đăng tải và giám sát chỉ số chuyển đổi theo thời gian thực.

Khi dữ liệu vận hành bị phân mảnh trên các bảng tính Excel rời rạc hoặc các công cụ giao tiếp phân tán, các rủi ro vận hành nghiêm trọng xuất hiện:
\begin{itemize}
 \item \textbf{Bất nhất thông tin}: Sai lệch số liệu giữa ngân sách kế hoạch và chi phí quảng cáo thực tế, thiếu tính truy vết về phiên bản nội dung cuối cùng đã được phê duyệt cũng như nguồn gốc tạo lập (\texttt{MANUAL} hay \texttt{AI}).
 \item \textbf{Lỗ hổng quy trình}: Thiếu cơ chế kiểm soát phân quyền dẫn đến nguy cơ nội dung chưa qua kiểm duyệt hoặc chứa từ khóa cấm (\texttt{banned\_keywords}) bị phát hành sai lệch tôn chỉ thương hiệu, hoặc người soạn thảo tự phê duyệt bài viết của chính mình.
 \item \textbf{Lãng phí nguồn lực sáng tạo}: Đội ngũ tiếp thị tiêu tốn phần lớn thời gian vào việc định dạng lại bài viết cho từng kênh hoặc tổng hợp báo cáo thủ công thay vì tập trung tối ưu chiến lược.
\end{itemize}

Sự bùng nổ của các mô hình ngôn ngữ lớn (Large Language Models -- LLMs) và kiến trúc tác tử kết nối công cụ (Model Context Protocol -- MCP) mở ra cơ hội tự động hóa sáng tạo nội dung và diễn giải dữ liệu. Tuy nhiên, nếu tích hợp LLM vào doanh nghiệp một cách tự phát và thiếu rào chắn kỹ thuật (\textit{guardrails}), hệ thống sẽ đối mặt với hiện tượng ảo giác (\textit{hallucination}), tự ý ghi đè dữ liệu nghiệp vụ hoặc không kiểm soát được trách nhiệm quyết định.

Đề tài \textbf{``Hệ thống quản lý chiến dịch marketing có tích hợp AI (MarketFlow AI)''} (Mã đề tài: 80300) được triển khai nhằm cung cấp giải pháp quản trị tập trung, kết hợp quy trình phê duyệt \textit{Human-in-the-Loop} (HITL) nghiêm ngặt với dịch vụ \texttt{ai-runner} cô lập và máy chủ nội bộ FastMCP kết nối 13 mô hình ngôn ngữ lớn đa nhà cung cấp. Hệ thống bảo đảm AI luôn đóng vai trò trợ lý đề xuất (\textit{Proposal-Only}), tuyệt đối không tự ý ghi đè cơ sở dữ liệu hay thay thế con người trong các quyết định kinh doanh nhạy cảm.

\section{Các bên liên quan và phân tích nhu cầu}

Hệ thống được thiết kế và hiện thực hóa nhằm phục vụ bốn nhóm tác nhân chính, với nhu cầu và tiêu chí bảo đảm cụ thể được lượng hóa trong Bảng~\ref{tab:stakeholders}.

\begin{table}[!htbp]
\centering
\footnotesize
\renewcommand{\arraystretch}{1.02}
\caption{Người sử dụng và tiêu chí bảo đảm kỹ thuật}
\label{tab:stakeholders}
\begin{tabularx}{\textwidth}{|L{3.2cm}|Y|Y|}
\hline
\thead{Bên liên quan} & \thead{Nhu cầu cốt lõi} & \thead{Điều hệ thống bảo đảm thực tế}\\
\hline
Người quản lý (\texttt{MANAGER}) & Quản trị vòng đời tài khoản nhân sự, thiết lập hồ sơ thương hiệu (\texttt{brand\_profile}), từ khóa cấm (\texttt{banned\_keywords}), danh mục 7 kênh; giám sát ngân sách, phê duyệt nội dung tiếp thị và theo dõi KPI (CTR, CVR, CPC, CPA, CPM). & Toàn quyền quản trị người dùng (\texttt{/api/v1/users}), kênh và Brand Kit; máy trạng thái HITL khóa chặt quyền phê duyệt (\texttt{/approve}, \texttt{/reject}) kèm quy tắc chống tự duyệt (\texttt{CANNOT\_REVIEW\_OWN}).\\
\hline
Nhân viên Marketing (\texttt{MARKETER}) & Khởi tạo chiến dịch được giao, gọi Trợ lý AI hỗ trợ lên ý tưởng/soạn bản nháp, chỉnh sửa nội dung, gửi duyệt, lập lịch và nhập chỉ số hiệu quả. & Phân quyền mức bản ghi (\texttt{owner\_id}): chỉ thao tác trên chiến dịch/nội dung thuộc sở hữu; giao diện SSE streaming trực quan với thẻ đề xuất \texttt{ai\_action = PROPOSED}.\\
\hline
Kỹ sư vận hành (DevOps/SecOps) & Đóng gói triển khai nhanh, cô lập tiến trình AI an toàn, dễ dàng sao lưu và kiểm thử tự động. & Dịch vụ \texttt{ai-runner} cô lập dưới tài khoản \texttt{mfai} (\texttt{MemoryMax=896M}), chặn \texttt{/mcp} từ Internet qua Nginx, CI/CD tự động kiểm thử 151 ca test.\\
\hline
\end{tabularx}
\end{table}

\section{Ranh giới và phạm vi hệ thống}

Hình~\ref{fig:context} thể hiện ranh giới hệ thống MarketFlow AI trong mối tương quan với các thực thể bên ngoài.

\begin{figure}[!htbp]
\centering
\vspace{2pt}
\resizebox{0.98\textwidth}{!}{%
\begin{tikzpicture}[
 font=\sffamily\footnotesize,
 >=Latex,
 user_node/.style={
 draw=black,
 fill=white,
 rounded corners=3pt,
 minimum width=3.2cm,
 minimum height=1.65cm,
 align=center,
 line width=0.85pt,
 text width=3.0cm,
 inner sep=3pt
 },
 sys_node/.style={
 draw=black,
 fill=black!4,
 rounded corners=4pt,
 minimum width=4.6cm,
 minimum height=2.7cm,
 align=center,
 line width=1.1pt,
 text width=4.3cm,
 inner sep=5pt
 },
 ext_node/.style={
 draw=black!80,
 fill=white,
 rounded corners=3pt,
 minimum width=3.2cm,
 minimum height=1.65cm,
 align=center,
 line width=0.85pt,
 dashed,
 text width=3.0cm,
 inner sep=3pt
 },
 rel/.style={-{Latex[length=2.2mm]},draw=black,line width=0.8pt},
 lab/.style={font=\sffamily\scriptsize,fill=white,inner sep=2pt,align=center,text=black}
]

\node[user_node] (mgr) at (0.8, 1.65) {
 \textbf{Người quản lý (MANAGER)}\\[2pt]
 {\scriptsize Quản trị kênh, Brand Profile, phê duyệt nội dung, xem báo cáo}
};

\node[user_node] (stf) at (0.8, -1.65) {
 \textbf{Nhân viên (MARKETER)}\\[2pt]
 {\scriptsize Quản lý chiến dịch, soạn bài, gọi Trợ lý AI, nhập \& import CSV số liệu}
};

\node[sys_node] (sys) at (7.0, 0) {
 \textbf{\small HỆ THỐNG MARKETFLOW}\\
 {\scriptsize (Đề tài AIA331 -- Mã số 80300)}\\[3pt]
 {\scriptsize Quản lý chiến dịch, 7 kênh, nội dung, duyệt bài HITL, lịch đăng, chỉ số KPI; tích hợp Trợ lý AI qua MCP \& cơ chế Chỉ Đề Xuất (PROPOSED)}
};

\node[ext_node] (ai) at (13.2, 1.65) {
 \textbf{Hạ tầng Đa mô hình AI}\\
 {\scriptsize (OpenCode Zen / Groq /}\\
 {\scriptsize NVIDIA NIM / Gemini)}\\[1pt]
 {\scriptsize Suy luận LLM \& gọi MCP Tool}
};

\node[ext_node] (soc) at (13.2, -1.65) {
 \textbf{7 Kênh Truyền thông}\\
 {\scriptsize (FB / TikTok / IG / Zalo /}\\
 {\scriptsize Email / Google Ads / Web)}\\[1pt]
 {\scriptsize Đăng bài đã duyệt \& lấy URL}
};

\draw[rel] (mgr.east) -- node[lab, pos=0.48]{Quản trị \&\\Phê duyệt} (sys.155);
\draw[rel] (stf.east) -- node[lab, pos=0.48]{Soạn thảo \&\\Xác nhận đề xuất} (sys.205);

\draw[rel] (sys.25) -- node[lab, pos=0.52]{Prompt /\\MCP Tool Calls} (ai.west);
\draw[rel] (sys.335) -- node[lab, pos=0.52]{Nội dung\\APPROVED} (soc.west);

\end{tikzpicture}%
}

\vspace{2pt}
\caption{Sơ đồ ranh giới hệ thống MarketFlow AI}
\label{fig:context}
\end{figure}

Trong phạm vi hiện thực hóa của hệ thống:
\begin{itemize}[noitemsep,topsep=2pt]
 \item \textbf{Bên trong ranh giới}: Ứng dụng Single Page Application (React 18 + TypeScript + Vite), máy chủ dịch vụ REST API kiêm máy chủ công cụ nội bộ FastMCP \texttt{/mcp} (FastAPI), cơ sở dữ liệu quan hệ 14 bảng (PostgreSQL 16 trên Production VPS và SQLite WAL mode khi phát triển/kiểm thử), cùng dịch vụ vi mô \texttt{ai-runner} (\texttt{127.0.0.1:8100}) điều phối tiến trình \texttt{opencode run}.
 \item \textbf{Ngoài ranh giới}: 5 cổng nhà cung cấp mô hình ngôn ngữ lớn (OpenCode Free, Groq, OpenRouter, Google AI Studio) kết nối qua HTTPS từ \texttt{ai-runner}, và các kênh phát hành nội dung bên ngoài (được quản lý thông qua quy trình duyệt nội bộ, lập lịch và xác nhận liên kết xuất bản \texttt{published\_url}).
\end{itemize}

\section{Mục tiêu và tiêu chí nghiệm thu thực chứng}

Hệ thống đặt ra 5 mục tiêu kỹ thuật cốt lõi (M1--M5), toàn bộ đều được hiện thực hóa và kiểm chứng bằng các bài kiểm thử tự động cũng như dữ liệu thực tế tại Bảng~\ref{tab:goals}.

\begin{table}[!htbp]
\centering
\footnotesize
\renewcommand{\arraystretch}{1.02}
\caption{Mục tiêu kỹ thuật và kết quả kiểm chứng thực tế}
\label{tab:goals}
\begin{tabularx}{\textwidth}{|C{1.0cm}|L{3.5cm}|Y|L{3.4cm}|}
\hline
\thead{Mã} & \thead{Mục tiêu} & \thead{Giải pháp kỹ thuật đã cài đặt} & \thead{Bằng chứng thực chứng}\\
\hline
M1 & Quản trị tập trung toàn bộ vòng đời tiếp thị trên 14 bảng CSDL. & CRUD hoàn chỉnh: Người dùng, Kênh, Chiến dịch, Nội dung, Lịch sử duyệt, Lịch đăng, Chỉ số (nhập tay + CSV 2 bước), Brand Kit, Thông báo. & Bộ kiểm thử Backend Pytest (104 ca kiểm thử trên 16 tệp test) đạt tỷ lệ Pass 100\%.\\
\hline
M2 & An toàn phân quyền RBAC + mức bản ghi và toàn vẹn quy trình HITL. & Phân quyền 2 vai trò (\texttt{MANAGER}, \texttt{MARKETER}) kết hợp kiểm tra \texttt{owner\_id}; quy tắc \texttt{CANNOT\_REVIEW\_OWN} và \texttt{CONTENT\_LOCKED} (HTTP 409) khi xóa bài đã duyệt. & Kiểm thử đối kháng bảo mật trên \texttt{test\_contents.py}, \texttt{test\_campaigns.py} chặn 100\% truy cập trái phép.\\
\hline
M3 & Tích hợp AI tác tử có kiểm soát (\textit{Proposal-Only}) và khử ảo giác. & Dịch vụ \texttt{ai-runner} + FastMCP 14 công cụ (9 Read, 5 Propose tạo \texttt{ai\_action = PROPOSED}); đối chiếu \texttt{check\_banned\_keywords}. & Đánh giá 4 giai đoạn tối ưu Prompt trên 8 kịch bản (bản Production đạt 7/8 -- 87,5\%); 18 ca test \texttt{ai-runner} + 38 ca test Assistant.\\
\hline
M4 & Trải nghiệm giao diện trực quan, truyền phát SSE thời gian thực. & Giao diện React 18 + Tailwind CSS v4, biểu đồ ECharts, lịch FullCalendar, khung chat SSE streaming hiển thị trực tiếp thẻ công cụ và đề xuất. & Kiểm thử TypeScript/ESLint 0 lỗi; 29 ca kiểm thử tự động đầu-cuối Playwright E2E trên Chromium đạt Pass 100\%.\\
\hline
M5 & Khả năng tái lập, đóng gói Docker và vận hành Production VPS. & \texttt{docker-compose.yml} 3 dịch vụ, kịch bản \texttt{seed\_demo.py}, triển khai thực tế trên VPS Ubuntu + PostgreSQL 16 + Systemd + Nginx. & Hệ thống vận hành ổn định tại \url{https://marketing.kienhieu.id.vn} với \texttt{health: ok, db: ok, ai\_runner: ok}.\\
\hline
\end{tabularx}
\end{table}

\section{Phạm vi hiện thực hóa}

\subsection{Các chức năng đã hiện thực hóa hoàn chỉnh}
\begin{enumerate}
 \item \textbf{Xác thực và quản trị nhân sự}: Đăng nhập và tự đăng ký JWT (\texttt{HS256}) kèm quản lý phiên \texttt{token\_version}, băm mật khẩu Bcrypt, kiểm soát 2 vai trò (\texttt{MANAGER}, \texttt{MARKETER}) và trạng thái hoạt động (\texttt{is\_active}).
 \item \textbf{Phân quyền mức bản ghi (Record-level Authorization)}: \texttt{MARKETER} chỉ được sửa/xóa chiến dịch, biên tập nội dung, lập lịch và nhập số liệu trên các chiến dịch do chính mình sở hữu (\texttt{owner\_id}). \texttt{MANAGER} có quyền quản trị và giám sát toàn cục.
 \item \textbf{Quản lý chiến dịch, kênh và hồ sơ thương hiệu}: Quản trị chiến dịch gắn với sản phẩm mục tiêu (\texttt{product}), đối tượng (\texttt{target\_audience}), ngân sách (\texttt{budget} $\ge 0$), thời gian (\texttt{start\_date} $\le$ \texttt{end\_date}), phân bổ ngân sách theo 7 loại kênh (\texttt{campaign\_channels}), cấu hình \texttt{brand\_profile} và từ khóa cấm (\texttt{banned\_keywords}).
 \item \textbf{Quy trình phê duyệt nội dung Human-in-the-loop}: Vòng đời trạng thái nghiêm ngặt: \texttt{DRAFT} $\rightarrow$ \texttt{PENDING\_REVIEW} (tự động quét từ khóa cấm khi \texttt{/submit}) $\rightarrow$ \texttt{APPROVED} / \texttt{REJECTED} (bởi \texttt{MANAGER}, chặn tự duyệt bài của chính mình qua mã lỗi \texttt{CANNOT\_REVIEW\_OWN}) $\rightarrow$ \texttt{PUBLISHED} (\texttt{/mark-published} kèm \texttt{published\_url}). Chặn xóa nội dung đã duyệt/xuất bản (\texttt{CONTENT\_LOCKED}) và tự động tạo thông báo (\texttt{notifications}).
 \item \textbf{Lập lịch xuất bản có điều kiện}: Tạo lịch đăng (\texttt{schedules}) đối với nội dung hợp lệ, đảm bảo thời gian lên lịch nằm trong khung thời gian của chiến dịch.
 \item \textbf{Ghi nhận chỉ số (nhập tay \& CSV 2 bước) và tính toán KPI}: Ghi nhận số liệu theo bộ ba duy nhất (\texttt{campaign\_id}, \texttt{channel\_id}, \texttt{metric\_date}); hỗ trợ tải tệp CSV xem trước (\texttt{/metrics/import/preview}) và xác nhận lưu (\texttt{/metrics/import/commit}); tự động tính toán CTR, CVR, CPC, CPA, CPM và tỷ lệ giải ngân ngân sách (\texttt{used\_pct}) với cơ chế chống chia cho 0.
 \item \textbf{Trợ lý AI tác tử qua FastMCP và SSE Streaming}: Hỗ trợ hội thoại nhiều luồng (\texttt{ai\_threads}, \texttt{ai\_messages}), lựa chọn trực tiếp 13 mô hình LLM đa nhà cung cấp trên giao diện, gọi 9 công cụ đọc dữ liệu thật và 5 công cụ tạo đề xuất (\texttt{ai\_actions} ở trạng thái \texttt{PROPOSED}) để người dùng xem xét bấm \textit{Xác nhận} (\texttt{EXECUTED}) hoặc \textit{Từ chối} (\texttt{REJECTED}).
\end{enumerate}

\subsection{Các nội dung chủ động loại trừ khỏi phạm vi}
Để tập trung tối đa nguồn lực vào độ tin cậy của kiến trúc phần mềm và kiểm soát an toàn AI, hệ thống chủ động đặt ngoài phạm vi:
\begin{itemize}
 \item Tự động trích tiền từ thẻ tín dụng hoặc kết nối trực tiếp vào cổng thanh toán ngân hàng.
 \item Tự động đẩy bài viết trực tiếp lên Graph API/TikTok API mà không qua xác nhận xuất bản của con người (sử dụng quy trình duyệt nội bộ và xác nhận \texttt{published\_url} tại bước \texttt{mark-published}).
 \item Tự huấn luyện (\textit{pre-training}) mô hình ngôn ngữ riêng từ đầu (sử dụng kỹ thuật điều phối tác tử MCP và tối ưu hóa System Prompt qua 4 giai đoạn cải tiến).
\end{itemize}

\section{Kết luận chương}
\enlargethispage{1.5\baselineskip}

Chương 1 đã định hình bối cảnh bài toán thực tế, xác lập các mục tiêu kỹ thuật định lượng và ranh giới hệ thống MarketFlow AI. Toàn bộ mục tiêu đề ra đều được quy chiếu vào các thành phần phần mềm cụ thể đã hoàn tất và vận hành thực tế. Chương tiếp theo sẽ phân tích cơ sở lý thuyết, các tiêu chuẩn kỹ thuật quốc tế và các quyết định công nghệ định hình nên kiến trúc hệ thống.
